% Part: first-order-logic
% Chapter: axiomatic-deduction
% Section: proof-theoretic-notions

\documentclass[../../../include/open-logic-section]{subfiles}

\begin{document}

\iftag{FOL}
      {\olfileid{fol}{axd}{ptn}}
      {\olfileid{pl}{axd}{ptn}}

\olsection{ثبوت-تيوريکي مفاهيم}

\begin{explain}
لکه څنګه چې مو يو شمېر مهم معنايي مفاهيم
(\iftag{FOL}{اعتبار}{صدق‌تابع قضيه}، معنايي لازميت او د صدق وړتيا)
تعريف کړي دي، اوس ورسره سم \emph{ثبوت-تيوريکي مفاهيم} تعريفوو۔
دا د !!{structure}s په دننه کښې د !!{sentence}s پوره کېدو ته په
رجوع نۀ، بلکې د ځينو فارمولونو !!{derivability} يا
!!{nonderivability} ته په رجوع تعريفېږي۔ دا يوه مهمه موندنه وه چې
دغه مفاهيم يو شان راځي۔ د هغو يو شان راتلل د \emph{سموالي} او
\emph{بشپړوالي د قضيو} منځپانګه ده۔
\end{explain}

\begin{defn}[!!^{derivability}]
!!^a{formula}~$!A$ هله له~$\Gamma$ څخه \emph{!!{derivable}} دے،
چې $\Gamma \Proves !A$ ليکو، چې له~$\Gamma$ څخه داسې
!!a{derivation} وي چې په~$!A$ پاے ته رسېږي۔
\end{defn}

\begin{defn}[قضيې]
!!^a{formula}~$!A$ هله \emph{قضيه} ده چې له تش سټ څخه د~$!A$ يو
!!a{derivation} وي۔ که $!A$ قضيه وي $\Proves !A$، او که نۀ وي
$\Proves/ !A$ ليکو۔
\end{defn}

\begin{defn}[سازګاري]
د !!{formula}s يو سټ~$\Gamma$ هله او يوازې هله \emph{سازګار} دے
چې $\Gamma\Proves/ \lfalse$؛ که داسې نۀ وي \emph{ناسازګار} دے۔
\end{defn}

\begin{prop}[انعکاس]
\ollabel{prop:reflexivity}
که $!A \in \Gamma$ وي، نو $\Gamma \Proves !A$۔
\end{prop}

\begin{proof}
  يوازې !!{formula}~$!A$ له~$\Gamma$ څخه د~$!A$ يو
  !!a{derivation} دے۔
\end{proof}

\begin{prop}[زياتېدنه]
\ollabel{prop:monotonicity}
که $\Gamma \subseteq \Delta$ او $\Gamma \Proves !A$ وي، نو
$\Delta \Proves !A$۔
\end{prop}

\begin{proof}
له~$\Gamma$ څخه د~$!A$ هر !!{derivation} له~$\Delta$ څخه هم د~$!A$
يو !!a{derivation} دے۔
\end{proof}

\begin{prop}[انتقال]
\ollabel{prop:transitivity}
که $\Gamma \Proves !A$ او $\{!A\} \cup \Delta \Proves !B$ وي،
نو $\Gamma \cup \Delta \Proves !B$۔
\end{prop}

\begin{proof}
  فرض کړئ $\{!A\} \cup \Delta \Proves !B$۔ بيا له
  $\{!A\} \cup \Delta$ څخه يو !!a{derivation}
  $!B_1$، \dots، $!B_l = !B$ شته۔ په دې !!{derivation} کښې ځينې
  ګامونه به د هغې قاعدې له مخې سم وي چې کومې مخکښې کرښې
  $!B_i = !A$ ته رجوع کوي۔ د فرض له مخې له~$\Gamma$ څخه د~$!A$
  يو !!a{derivation} شته؛ يعنې داسې !!a{derivation}
  $!A_1$، \dots، $!A_k = !A$ چې هر $!A_i$ يا بديهي اصل وي، يا
  د~$\Gamma$ يو !!a{element} وي، يا د استنتاج د قاعدې له مخې سم
  وي۔ اوس دا لړۍ په پام کښې ونيسئ:
  \[
  !A_1, \dots, !A_k = !A, !B_1, \dots, !B_l = !B.
  \]
  دا له~$\Gamma \cup \Delta$ څخه د~$!B$ سم !!{derivation} دے؛ ځکه
  هر $!B_i = !A$ اوس د هماغې قاعدې له مخې توجيه کېږي چې
  $!A_k = !A$ توجيه کوي۔
  \textbf{د ژباړې د نسخې سمون (OLAX-002):} سرچينې په وروستۍ جمله
  کښې يوازې له بي-آی څخه د ميتا-فارمول نښه پرېښې وه؛ هماغه کرښه
  پورته او په نښلول شوې لړۍ کښې په نښه لرلې بڼه تعريف شوې ده۔
\end{proof}

پام وکړئ چې په ځانګړي ډول دا معنٰی لري: که
$\Gamma \Proves !A$ او $!A \Proves !B$ وي، نو
$\Gamma \Proves !B$۔ له دې څخه دا هم راځي چې که
$!A_1, \dots, !A_n \Proves !B$ او د هر~$i$ دپاره
$\Gamma \Proves !A_i$ وي، نو $\Gamma \Proves !B$۔

\begin{prop}
\ollabel{prop:incons}
$\Gamma$ هله او يوازې هله ناسازګار دے چې د هر~$!A$ دپاره
$\Gamma \Proves !A$ وي۔
\end{prop}

\begin{proof}
تمرين۔
\end{proof}

\tagprob{FOL}
\begin{prob}
\olref[fol][axd][ptn]{prop:incons} ثابت کړئ۔
\end{prob}
\tagendprob

\tagprob{notFOL}
\begin{prob}
\olref[pl][axd][ptn]{prop:incons} ثابت کړئ۔
\end{prob}
\tagendprob

\begin{prop}[تراکم]
\ollabel{prop:proves-compact}
  \begin{enumerate}
  \item که $\Gamma \Proves !A$ وي، نو داسې متناهي فرعي سټ
    $\Gamma_0 \subseteq \Gamma$ شته چې $\Gamma_0 \Proves !A$۔
  \item که د~$\Gamma$ هر متناهي فرعي سټ سازګار وي، نو $\Gamma$
    سازګار دے۔
  \end{enumerate}
\end{prop}

\begin{proof}
  \begin{enumerate}
    \item که $\Gamma \Proves !A$ وي، نو د !!{formula}s داسې متناهي
      لړۍ $!A_1$، \dots،~$!A_n$ شته چې $!A \ident !A_n$؛ او هر
      $!A_i$ يا منطقي بديهي اصل وي، يا د~$\Gamma$ يو !!a{element}
      وي، يا د مودوس پوننس له مخې له مخکښو !!{formula}s څخه راځي۔
      $\Gamma_0$ د هغو $!A_i$ سټ ونيسئ چې په~$\Gamma$ کښې دي۔ بيا
      همدا !!{derivation} له~$\Gamma_0$ څخه هم يو !!a{derivation}
      دے، نو $\Gamma_0 \Proves !A$۔
    \item دا د ځانګړي حالت~$!A \ident \lfalse$ دپاره د~(۱) مقابل
      عکس دے۔
\end{enumerate}
\end{proof}

\end{document}
